\documentclass[11pt]{article}
\usepackage{shared/template_tgir_article}
\usepackage{booktabs}
\usepackage{longtable}
\usepackage{tabularx}
\usepackage{array}
\usepackage{enumitem}
\usepackage{xcolor}
\usepackage{listings}

\lstset{
    basicstyle=\ttfamily\small,
    breaklines=true,
    columns=fullflexible,
    frame=single
}

\newcommand{\product}{TGIR QuantLib Tools}
\newcommand{\code}[1]{\texttt{#1}}
\newcommand{\docdate}{2026-04-06}

% Visual design is controlled exclusively by template_tgir_article.sty and
% IEEEtran. Keep this file limited to content-support packages and commands.


\title{IT Operations Guide\\\large \product}
\author{TG Investments and Research}
\date{\docdate}

\begin{document}
\maketitle
\tableofcontents
\newpage

\section{Purpose}
This guide is for IT and operations staff responsible for runtime support, access control, incident response, backups, and production hygiene.

\section{Runtime Summary}
The application is a single Flask service with no database. Operational state is limited to:
\begin{itemize}[leftmargin=1.5em]
    \item the deployed code,
    \item the Python virtual environment,
    \item the local \code{.env} file,
    \item logs and service definitions,
    \item the generated debug CSV under the app's debug directory, and
    \item standard operating-system service state.
\end{itemize}

\section{Minimum Environment Requirements}
\begin{itemize}[leftmargin=1.5em]
    \item Linux host or local macOS workstation.
    \item Python 3 with the ability to create \code{.venv/}.
    \item QuantLib Python package installation support.
    \item Apache reverse proxy and a local WSGI service for production.
    \item Read-write access to the debug output path.
\end{itemize}

\section{Secrets And Configuration}
\subsection{Required Variables}
\begin{tabularx}{\textwidth}{>{\raggedright\arraybackslash}p{0.32\textwidth}X}
\toprule
Variable & Purpose \\
\midrule
\code{FLASK\_SECRET\_KEY} & Signs Flask sessions. \\
\code{APP\_LOGIN\_USERNAME} & Single application login username. \\
\code{APP\_LOGIN\_PASSWORD} or \code{APP\_LOGIN\_PASSWORD\_HASH} & Login credential. \\
\code{PORT} or \code{FLASK\_RUN\_PORT} & Local port override if required. \\
\bottomrule
\end{tabularx}

\subsection{Storage Rules}
Do not commit secrets. Store production secrets in a server-side \code{.env} file outside Git-controlled directories, and use least-privilege filesystem permissions.

\section{Operational Tasks}
\subsection{Start And Stop}
Local development uses \code{./.venv/bin/python app.py}. Production should use a systemd-managed WSGI service rather than the Flask development server.

\subsection{Health Check}
Use \code{/health} as the first smoke check after restart or deployment. It should return a JSON document with \code{"ok": true}.

\subsection{Logs}
Primary production logs should come from:
\begin{itemize}[leftmargin=1.5em]
    \item systemd journal for the WSGI process,
    \item Apache access and error logs, and
    \item CI/CD workflow logs for deployment history.
\end{itemize}

\section{Security Checklist}
\begin{itemize}[leftmargin=1.5em]
    \item Keep \code{FLASK\_SECRET\_KEY} strong and unique.
    \item Run behind HTTPS in production.
    \item Restrict shell access to the droplet.
    \item Patch the host OS and Apache packages regularly.
    \item Use a dedicated service account when practical.
    \item Keep the \code{.env} file outside world-readable paths.
\end{itemize}

\section{Backup And Recovery}
There is no database to back up. Recovery focuses on:
\begin{itemize}[leftmargin=1.5em]
    \item code releases,
    \item the production \code{.env} file,
    \item Apache and systemd unit templates, and
    \item any shared debug or log directories kept outside the current release.
\end{itemize}
The deployment guide uses a \code{current} symlink plus timestamped release directories so rollback can be a symlink change followed by a service restart.

\section{Troubleshooting}
\begin{tabularx}{\textwidth}{>{\raggedright\arraybackslash}p{0.28\textwidth}X}
\toprule
Symptom & Operational response \\
\midrule
Service will not start & Confirm Python package install, \code{.env} presence, and systemd working directory. \\
\code{/health} fails & Check local bind port, WSGI service, and Apache proxy settings. \\
Login fails in production & Verify the deployed \code{.env} credential values and file permissions. \\
Disk growth & Clear old release directories, old logs, and generated artifacts according to retention policy. \\
\bottomrule
\end{tabularx}

\section{Reference Files}
Operations should know these files:
\begin{itemize}[leftmargin=1.5em]
    \item \path{wsgi.py},
    \item \path{requirements-production.txt},
    \item \path{deploy/systemd/tgir-quantlib-tools.service.template},
    \item \path{deploy/apache/tgir-quantlib-tools.conf.template}, and
    \item \path{.github/workflows/deploy_digitalocean.yml}.
\end{itemize}

\end{document}
